
\begin{tikzpicture}[x=1cm,y=1cm,every node/.style={font=\small,align=center},line width=.45pt]
\node (a) at (0,0) {先前行为制造的危险};\node (b) at (5.2,0) {介入因素制造的危险};\node (c) at (10,0) {实害结果};
\draw[->] (a)--(b);\draw[->] (b)--(c);
\end{tikzpicture}

% 人工对照底本逐节点转写；独立TikZ重绘。

\begin{tikzpicture}[x=1cm,y=-1cm,every node/.style={anchor=west,align=left,inner sep=2pt,font=\fontsize{9.5}{12.3}\selectfont},edge/.style={draw=black!65,line width=.45pt}]
\node[text width=2.200cm] (n1) at (0.000,1.446) {介入因素是否异常};
\node[text width=2.791cm] (n2) at (2.620,0.275) {不异常：引发关系};
\node[text width=5.128cm] (n3) at (5.840,0.275) {结论：先前行为与结果有因果关系};
\node[text width=2.457cm] (n4) at (2.620,1.837) {异常：独立关系};
\node[text width=2.700cm] (n5) at (5.840,1.837) {第二步：谁的作用大};
\node[text width=7.040cm] (n6) at (8.960,1.056) {先前行为的作用大：先前行为与结果有因果关系};
\node[text width=7.040cm] (n7) at (8.960,1.837) {介入因素的作用大：介入因素与结果有因果关系};
\node[text width=3.793cm] (n8) at (8.960,2.618) {二者作用都大：二因一果};
\draw[edge,->] (n2.east)--(5.690,0.275)--(5.690,0.275)--(n3.west);
\draw[edge] (n1.east)--(2.470,1.446)--(2.470,0.275)--(n2.west);
\draw[edge] (n5.east)--(8.810,1.837)--(8.810,1.056)--(n6.west);
\draw[edge] (n5.east)--(8.810,1.837)--(8.810,1.837)--(n7.west);
\draw[edge] (n5.east)--(8.810,1.837)--(8.810,2.618)--(n8.west);
\draw[edge,->] (n4.east)--(5.690,1.837)--(5.690,1.837)--(n5.west);
\draw[edge] (n1.east)--(2.470,1.446)--(2.470,1.837)--(n4.west);
\end{tikzpicture}

